\capitulo{4}{Técnicas y herramientas}

\section{Metodología de desarrollo}

El proyecto se ha desarrollado siguiendo principios de metodologías ágiles utilizando iteraciones organizadas en sprints.

Para la gestión de tareas se ha utilizado Zube.io, permitiendo organizar historias de usuario, tareas técnicas y seguimiento del progreso.

\section{Lenguaje y framework backend}

La aplicación backend se ha desarrollado utilizando Python junto con el framework Flask.

Flask permite construir APIs REST ligeras y fáciles de mantener.

\section{Control de versiones}

El control de versiones se ha realizado utilizando Git y GitHub.

Se han realizado commits frecuentes para mantener trazabilidad completa del desarrollo.

\section{Plataforma cloud}

La infraestructura cloud se basa en Microsoft Azure.

Los servicios principales utilizados son:

\begin{itemize}
    \item Azure App Service.
    \item Azure Storage.
    \item Azure Key Vault.
    \item Azure Monitor.
\end{itemize}

\section{Integración continua}

Se ha configurado un pipeline CI/CD utilizando GitHub Actions.

Este sistema automatiza:

\begin{itemize}
    \item Validación del código.
    \item Preparación del despliegue.
    \item Automatización de tareas.
\end{itemize}

\section{Calidad del código}

La calidad del código se analiza mediante SonarQube.

Esta herramienta permite detectar:

\begin{itemize}
    \item Vulnerabilidades.
    \item Duplicidad.
    \item Complejidad excesiva.
    \item Posibles errores.
\end{itemize}